% ===================== CONFIGURACIÓN DE PÁGINA MEJORADA =====================
\usepackage[top=4.5cm, bottom=2.5cm, left=2.0cm, right=2.0cm,
headheight=4cm, headsep=0.5cm]{geometry}

% ===================== TIPOGRAFÍA Y CODIFICACIÓN =====================
\usepackage[utf8]{inputenc}
\usepackage[T1]{fontenc}
\usepackage[spanish,es-tabla]{babel}
\usepackage{times}
\usepackage{microtype}

% ===================== ELIMINAR SANGRÍAS =====================
\setlength{\parindent}{0pt}
\setlength{\parskip}{6pt}

% ===================== CONFIGURACIÓN DE TÍTULOS =====================
\usepackage{titlesec}
\titleformat{\section}{\bfseries\fontsize{12}{14}\selectfont}{\thesection}{1em}{}
\titleformat{\subsection}{\bfseries\fontsize{12}{14}\selectfont}{\thesubsection}{1em}{}
\titleformat{\subsubsection}{\bfseries\fontsize{12}{14}\selectfont}{\thesubsubsection}{1em}{}

% ===================== PAQUETES MATEMÁTICOS =====================
\usepackage{amsmath,amssymb,amsfonts,mathtools}
\usepackage{textgreek,gensymb,siunitx}

% ===================== QUÍMICA Y SÍMBOLOS =====================
\usepackage[version=4]{mhchem}
\usepackage{chemfig,textcomp}

% ===================== TABLAS Y FIGURAS =====================
\usepackage{booktabs,array,longtable,multirow}
\usepackage{graphicx,float,subcaption}
\usepackage{tabularx}

% ===================== SOPORTE PARA SVG =====================
\usepackage{svg}

% ===================== ENLACES E HIPERVÍNCULOS =====================
\usepackage[hidelinks,colorlinks=true,linkcolor=blue,citecolor=blue]{hyperref}
\usepackage{url}

% ===================== OTROS PAQUETES =====================
\usepackage{enumitem,xcolor,listings,algorithm,algpseudocode}
\usepackage{tocloft}
\usepackage{appendix}

% ===================== PAQUETES ADICIONALES PARA TABLAS ESTÉTICAS =====================
\usepackage{array}
\usepackage{colortbl}
\usepackage{multirow}
\usepackage{booktabs}

% ===================== CONFIGURACIÓN DEL ENCABEZADO =====================
\usepackage{fancyhdr}

% ===================== INFORMACIÓN DEL DOCUMENTO =====================
\newcommand{\documentcode}{P25XX-PR-INF-00X}
\newcommand{\targetcompany}{Empresa Cliente}
\newcommand{\projecttitle}{Análisis Técnico de Sistema de Ingeniería}

% ===================== CONFIGURACIÓN DE LISTINGS =====================
\lstset{
	basicstyle=\footnotesize\ttfamily,
	breaklines=true,
	frame=single,
	language=Python,
	showstringspaces=false
}

\journal{Revista de Ingeniería}